% ===== 原文 pp.63–66（PDF p.73–76） =====
% [原文 p.63]
从 3d 视角看，这意味着相应的边界线算子是 $\cA$ 在表示 $V$ 中的 Wilson 线算子。在经典层面上，表示 $V_1$ 与 $V_2$ 中两个 Wilson 线算子的融合给出表示 $V_1\otimes V_2$ 中的 Wilson 线，并且显然这一结果不会有量子修正（规范耦合 $e^2$ 是无关参数）。

% [Tutor Note: p.63 "无量子修正"的核对：3d 规范耦合 $e^2$ 的质量量纲为 1，是无关耦合，故 Wilson 线的融合只由经典表示论（$V_1\otimes V_2$）决定；这与 p.59 说 Gukov–Witten 面算子对任意 $t$ 的 BRST 不变性互不矛盾（那里 $t$ 是扭转参数，不是耦合）。]

总结一下，幺半范畴 $\mathbb{V}_{\mathbb{X}_0\mathbb{X}_0}$ 就是 $\C^*$ 的有限维分次表示范畴，或者等价地，是单点上凝聚层的等变导出范畴，我们把它记作 $D^b_{\C^*}(Coh(\bullet))$。我们提出：边界条件的 2-范畴等价于 $D^b_{\C^*}(Coh(\bullet))$ 上模范畴所构成的 2-范畴。为给出这类模范畴的一类具体例子，考虑一个带 $\C^*$ 作用的 Calabi-Yau 流形 $Y$。相应的 B-模型可以耦合到边界规范场，从而为 3d 规范理论提供一组自然的拓扑边界自由度。相应的边界改变线算子范畴是 $Y$ 的 $\C^*$-等变有界导出范畴，它显然是一个 $D^b_{\C^*}(Coh(\bullet))$ 上的模范畴。

\subsection{把各个扇区拼起来}

如何把两个模型的边界条件组合起来是相当显然的。完整理论中最基本的边界条件是：RW 扇区中 $\sigma=0$，规范扇区中自由（Neumann）条件。我们把它称为有特殊地位的边界条件。在边界上自由的玻色场是 $\C^*$ 值的标量 $h=\exp(-2\pi\tau)$，以及复规范场 $\cA=A+\ii\phi$ 的限制。更一般的边界条件涉及纤维化在 $\C^*$ 上、并容许一个 $\C^*$ 作用的边界 B-模型或边界 Landau-Ginzburg 模型。在 $\C^*$ 上的纤维化确定了与 $\tau$ 的边界值的耦合，而 $\C^*$ 作用确定了与边界规范场 $\cA$ 的耦合。

与 RW 模型中一样，我们可以对完整理论中所有边界条件的集合给出一个更代数的定义。这个描述也有用，因为它提示了如何定义边界条件集合的 2-范畴结构。我们考虑有特殊地位的那个边界条件所对应的边界线算子幺半范畴。它就是由带规范场的 RW 模型在区间上约化所得的 2d TFT 的膜范畴。由于 B-型 3d 规范理论的约化给出 B-型 2d 规范理论，而 RW 模型的约化给出以 $\C^*$ 为靶空间的 B-模型，所以有效的 2d TFT 就是以 $\C^*$ 为靶空间的带规范场 B-模型，其中规范群 $U(1)$ 平凡地作用。如 §A.3 所述，相应的膜范畴等价于 $D^b_{\C^*}(Coh(\C^*))$。幺半结构无法由 2d 考虑确定，但容易看出（给定 RW 模型与 3d B-型规范理论的结果）它由导出张量积给出。

% [Tutor Note: p.63 标量 $h=\exp(-2\pi\tau)$ 的取值核对：按 §2 的约定 $\tau=\theta/2\pi+4\pi\ii/e^2$，故 $h=\exp(-\ii\theta-8\pi^2/e^2)$，模长 $|h|=\exp(-8\pi^2/e^2)>0$，确为 $\C^*$ 值；而 p.64 中 $h_0=\exp(-2\pi(\alpha-\ii\gamma))$ 与 p.59 奇点沿小圆的 $\cA$ 和乐 $\exp(2\pi(\alpha-\ii\gamma))$ 互为逆（同一约定下），比对符号时以此为准。]

% [原文 p.64]
每个边界条件都给出该幺半范畴上的一个模范畴。自然可以猜测其逆也成立，即：该幺半范畴上每个合理的模范畴都可以看作完整 3d TFT 的一个边界条件。例如，我们可以考虑由 $\C^*$ 的点所参数化的一族 Calabi-Yau 流形，使得族中每个模型都有 $\C^*$ 对称性。相应的模范畴就是该纤维化全空间的 $\C^*$-等变导出范畴。这就给出了母 4d 规范理论中面算子 2-范畴的一个猜测性描述。

% [Tutor Note: p.64 "合理的模范畴"这一限制条件不可省：$D^b_{\C^*}(Coh(\C^*))$ 上的模范畴在导出意义下未必都来自几何（可参见 [36,37] 与 Toen–Vezzosi [44] 的框架）；此处只是猜测（conjectural），原文用 "we propose / natural to conjecture"，翻译时保留其猜测地位，不要写成定理。]

下面描述 Gukov-Witten 面算子如何纳入这一图景。这类算子依赖于一个复参数 $h_0=\exp(-2\pi(\alpha-\ii\gamma))$，取值于 $\C^*$。从 3d 视角看，$h_0$ 确定了 RW 扇区中标量 $h=\exp(-2\pi\tau)$ 的边界值。另一个标量 $\sigma$ 则保持自由。因此 RW 扇区的边界条件对应于 $T^*[2]\C^*$ 中由 $h=h_0$ 给出的拉格朗日子流形（即 $h_0$ 点上的纤维）。规范扇区的边界条件是 Neumann 型的，没有非平凡参数。这里没有边界自由度。从我们的代数视角可以这样描述：在通常的 RW 理论中，$h=h_0$ 上的纤维对应于 $\C^*$ 上一个 DG-范畴的摩天层，它在 $h_0$ 上的"茎"是有界向量空间复形的范畴。我们可以把它记作 $D^b(Coh(\bullet))$。把规范自由度包括进来，意味着处理一个带 $\C^*$ 作用的范畴层。于是我们只需考虑 $\C^*$ 上一个范畴的摩天层，它在 $h_0$ 上的"茎"是 $\C^*$-等变向量空间复形的范畴 $D^b_{\C^*}(Coh(\bullet))$。幺半范畴 $D^b_{\C^*}(Coh(\C^*))$ 以相当显然的方式作用在它上面：只需把 $D^b_{\C^*}(Coh(\bullet))$ 的一个对象与 $D^b_{\C^*}(Coh(\C^*))$ 的一个对象到点 $h=h_0$ 的（导出）限制作张量。

% [Tutor Note: p.64 原文此处有排印问题："have no nontrivial parameters.There are no boundary degrees of freedom."（句号后缺空格），以及 "$D^b_{\C^*}\left(Coh(\C^*)\right)$ 的 $C^\ast$ -action"等处多余空格；译文按正常中文排版处理，无需改动内容。]

\subsection{Gukov-Witten 面算子上的线算子}

两个不同范畴的摩天层之间的态射范畴是平凡的（对象集为空）。这对应于一个事实：两个不同的 Gukov-Witten 面算子不能沿一条边界改变线算子相接。但是，位于某个特定 Gukov-Witten 面算子上的线算子范畴（即 Gukov-Witten 面算子的自同态范畴）是非平凡的。它最明显的对象是复化规范场 $\cA$ 的 Wilson 线，它们显然是 BRST 不变的。这类算子由 $\C^*$ 的不可约表示标记。因此人们或许会猜测：面算子上的线算子范畴就是 $\C^*$ 的表示范畴，或者是我们上面记作 $D^b_{\C^*}(Coh(\bullet))$ 的 $\C^*$-等变向量空间复形的范畴。然而，这一朴素猜测是错误的，可以通过考察能插入到这种 Wilson 线算子中的 BRST 不变局部算子看出来。从抽象视角看，它们构成一个代数（线算子范畴中某个对象的自同态代数）。显然场 $\sigma$ 的任意幂都给出这样的算子，所以线算子上的局部算子代数就是单变量、鬼数为 2 的多项式函数代数。在下文中我们把由 $\sigma$ 参数化的那条线记作 $\C[2]$，以表明

% [原文 p.65]
$\sigma$ 处于次数 2；因此 $\C[2]$ 是一个纯偶的分次流形。另一方面，$\C^*$ 的不可约表示的自同态代数就是 $\C$。

% [Tutor Note: p.64–65 自洽性核对：$\C[\sigma]$ 的 $\sigma$ 的鬼数为 2，故 $\C[2]$ 的坐标环是偶的；而 $\C^*$ 的不可约表示（一维、特征标 $n$）的自同态代数是 $\C$（Schur 引理）。下文 p.65 的结论"对象是 $\C[2]$ 上等变全纯向量丛复形，自同态是 $\C[\sigma]$"正是这两点的结合。]

要确定线算子范畴是什么，方便的做法是采取 2d 视角，把 3d 理论在带 Gukov-Witten 型边界条件（两端都是）的区间上约化。设 $x^3$ 是区间上的坐标。Gukov-Witten 边界条件消去了复标量 $h$（此时它被锁定在值 $h_0$）与场 $\phi_3$，但保留复标量 $\sigma$ 与规范场 $\cA$。因此有效 2d 理论也有两个扇区：以 $\C[2]$ 为靶空间的 B-模型，以及 B-型 2d 拓扑规范理论。根据 A.3，相应的膜范畴等价于 $\C[2]$ 的 $\C^*$-等变导出范畴：其对象可以看作 $\C[2]$ 上等变全纯向量丛的 $\C^*$-等变复形（$\C[2]$ 上的 $\C^*$ 作用平凡）。这个答案不依赖于 Gukov-Witten 面算子的参数 $h_0=\exp(-2\pi(\alpha-\ii\gamma))$。特别地，我们可以取平凡的面算子 $h_0=1$，此时应当得到 $t=\ii$ 时 GL-扭转理论中的体线算子范畴。

不难看出，这个关于体线算子范畴的答案与上面解释的 Wilson 线自同态代数的计算一致。事实上，插入一条 Wilson 线不对 $\sigma$ 施加任何约束，也不增加任何自由度，因此应当对应于 $\C[2]$ 上的平凡线丛。其纤维携带由 Wilson 线电荷确定的 $\C^*$ 表示。$D^b_{\C^*}(Coh(\C[2]))$ 中这类对象的自同态代数就是 $\C[2]$ 上的多项式函数代数。

% [Tutor Note: p.65 这个"平凡线丛"的对应可这样验证：$\C[2]$ 上平凡线丛 $\cO$ 的截面环是 $\C[\sigma]$，而 Wilson 线插入不改变 $\sigma$ 的取值空间（不增自由度），故其自同态 = $\C[\sigma]$，与 p.64 从 Wilson 线直接读出的局部算子代数相符；电荷 $n$ 只体现在 $\C^*$ 等变结构（权 $n$）上。]

现在清楚了：线算子范畴包含 Wilson 线之外的对象。例如，我们可以考虑 $\C[2]$ 原点处的摩天层，其在原点的茎是携带某个 $\C^*$ 表示的一维复线 $V$。定义相应的线算子有两种不同方式。第一种，我们可以考虑该摩天层的一个自由消解：
\[
V[-2]\otimes\cO\to V\otimes\cO ,
\]
其中 $V[-2]$ 表示把 $V$ 放在鬼数 $-2$，$\cO$ 是 $\C[2]$ 上的多项式函数代数，而上链映射是乘以 $\sigma$。此处平移 $-2$ 是必要的，以便上链映射的总次数为 1。这样一个消解的存在意味着，我们可以把"摩天层"线算子实现为两条 Wilson 线的"束缚态"，这两条 Wilson 线都联系着表示 $V$，但被放在不同的上同调次数上。相应的体线算子用 §A.3 的公式得到，其中带规范场 B-模型的靶空间取为 $\C[2]$，$\C[2]$ 上的向量丛 $E$ 平凡且秩为 2，其分次分量位于次数 1 与 0，而从前者到后者的丛态射 $T$ 是乘以 $\sigma$。与 §A.3 的

% [Tutor Note: p.65 次数记账需复核：按 $\sigma$ 的鬼数为 2（p.65 开头已定），
% 在 $V[-2]\otimes\cO\to V\otimes\cO$ 中"乘以 $\sigma$"把次数 $-2$ 的元素送到次数 0，映射的净次数为 0，并非原文所说"总次数 1"；
% 而 p.66 又要求 $T$ 把次数 1 的分量映到次数 0 的分量。这两处与 §A.3 的超联络公式（其非对角元 $\frac12(\psi-\ii\tilde\psi)$ 的次数为 1）需要统一复算，此处存疑，原文照译未改。]

% [原文 p.66]
结果一致，我们考虑 $\sigma^*E$ 上形如
\[
\cN=\begin{pmatrix}
n\cA & 0\\[2pt]
\frac{1}{2}(\psi-\ii\tilde\psi) & n\cA
\end{pmatrix}
\]
的超联络，其中 $n\in\Z$ 是 $\C^*$ 作用在 $V$ 上的权。对应于 $\C[2]$ 原点处摩天层的体线算子，就是这个超联络沿插入线 $\ell$ 的和乐。

另一种（等价的）方式是认真对待"摩天层定域在 $\sigma=0$ 处"这一事实，并要求场 $\sigma$ 在插入线 $\ell$ 上为零。为使这一点有良好定义，需要挖去 $\ell$ 的一个小管状邻域，并在所得边界上施加合适的边界条件。这个条件必须令 $\sigma=0$，同时不约束 $\cA$ 与边界 $\ell$ 相切的分量。BRST 不变性唯一地确定了所有其他场的边界条件。

% [Tutor Note: p.66 超联络矩阵的核对：对角元 $n\cA$ 中的 $n\in\Z$ 是 $\C^*$ 权，故取不同 $n$ 给出不同的线算子（对应 $V$ 的不同电荷）；非对角元 $\frac12(\psi-\ii\tilde\psi)$ 是 1-形式部分，正是 §A.3 中"丛态射 $T$"的场论实现。$\C[2]$ 上的 $\C^*$ 作用平凡这一事实保证了 $n$ 只能来自 $V$ 的权。]

\section{\texorpdfstring{$t=\ii$}{t=i} 且 \texorpdfstring{$G$}{G} 非阿贝尔时的面算子}

§5.5 已表明：对一般 $G$，$t=\ii$ 理论的紧化给出以 $T^*[2]G_\C$ 为靶空间的带规范场 Rozansky-Witten 模型（"[2]"表示余切丛的纤维坐标被形式地赋予鬼数 2）。$T^*[2]G_\C$ 上的 $G$-作用由底空间上的共轭作用诱导。

% [Tutor Note: p.66 非阿贝尔情形与 §6.2 的衔接：$G=U(1)$ 时 $T^*[2]G_\C=T^*[2]\C^*$，共轭作用平凡，与 §6.2 的所有结果一致；非阿贝尔时底空间的 $G_\C$-共轭作用使边界条件的分类变成"$G$-不变的复拉格朗日子流形"的分类，这是 §6.3.1 的全部内容。核对时可用 $G=U(1)$ 退化检验每条陈述。]

% [Tutor Note: p.66 余切丛上的 $G$-作用：底空间 $G_\C$ 上的共轭作用 $g\mapsto hgh^{-1}$ 提升到 $T^*G_\C$ 后，在纤维（$\cong\fg_\C^*$）上给出的是余伴随作用；这与 $\sigma$ 取值于 $\fg_\C$、按伴随表示变换是一致的（用不变双线性型把 $\fg_\C$ 与 $\fg_\C^*$ 等同）。]

\subsection{一些简单的边界条件}

我们来看以 $T^*[2]G_\C$ 为靶空间的带规范场 Rozansky-Witten 模型中的一些边界条件。规范扇区中最自然的边界条件是 Neumann 条件，它在边界上保持完全的规范不变性。在物质扇区中，必须选取 $T^*[2]G_\C$ 的一个 $G$-不变复拉格朗日子流形，且它要对纤维的伸缩不变。这样的拉格朗日子流形可以这样构造：取 $G_\C$ 的一个 $G_\C$-不变闭复子流形，再取它的余法丛。例如可以取整个 $G_\C$，此时拉格朗日子流形由 $\sigma=0$ 给出。我们把带规范场 RW 模型中由此得到的边界条件称为有特殊地位的边界条件。它是阿贝尔情形中 NN 条件的类比。

% [Tutor Note: p.66 余法丛是复拉格朗日的定维核对：$T^*G_\C$ 的复维为 $2n$（$n=\dim_\C G_\C$）；复子流形 $S$ 的余法丛的复维为 $\dim_\C S+(n-\dim_\C S)=n$，恰为全空间复维的一半，故对任何复子流形都自动是复拉格朗日（且对纤维伸缩不变）。取 $S=G_\C$ 得零截面 $\sigma=0$；取 $S$ 为共轭类得本小节第二段的构造。]

另一个自然的 $G$-不变拉格朗日子流形选择，是 $G_\C$ 中某个复共轭类的余法丛。为使该子流形是闭的，需取共轭类为半单的。这个边界条件是 ND 条件的非阿贝尔类比。相应的面算子是半单的 Gukov-Witten 型面算子。事实上，固定 $\exp(-2\pi(A_4+\ii\phi_4))$ 的一个半单共轭类，等同于固定 4d 规范理论中复联络 $A+\ii\phi$ 的极限和乐的一个半单共轭类。

% [Tutor Note: p.66 "NN"与"ND"两个标签的约定核对：第一个字母指规范扇区（N = Neumann），第二个字母指物质扇区（N = 自由/Neumann，D = 固定和乐/Dirichlet）。故"整个 $G_\C$ 的余法丛 $\sigma=0$"（物质扇区不固定任何量）是 NN 的类比，而"共轭类的余法丛"（固定了极限和乐所在的共轭类）是 ND 的类比。此命名与 §6.4 中规范扇区的 D/N 用法一致。]
